% ============================================================================
%  PHỤ LỤC
% ============================================================================
\clearpage
\appendix
\section*{Phụ lục}
\addcontentsline{toc}{section}{Phụ lục}
\renewcommand{\thesubsection}{\Alph{subsection}}

\subsection{Quy trình thực nghiệm bắt buộc}
\label{app:protocol}

Phụ lục này là quy trình nhóm thực hiện phải theo khi chạy E0--E5. Lệnh chính xác, đủ cờ, nằm trong
\code{code/RUNBOOK.md}; khi phụ lục và RUNBOOK khác nhau về lệnh, theo RUNBOOK. Khi khác nhau về phương pháp, theo
đề cương (\code{docs/design/2026-11-12-lung-nodule-cxr-vlm-lora-design.md}, mục 8.5). Năm quy tắc:
\begin{enumerate}
  \item Chạy theo đúng thứ tự ở Hình~\ref{fig:workflow}. Không chạy E1--E5 khi còn mục kiểm tra trước chưa đạt.
  \item Chỉ chọn ngưỡng và cấu hình trên validation, tức là từ hai file \code{thresholds.json} và
        \code{val\_selection.csv}. Không mở cột test để chọn, kể cả khi số trên test ``đẹp'' hơn.
  \item Không gõ tay số vào báo cáo. Mọi ô \textcolor{badred}{[--]} chép từ file trong \code{outputs/} như
        Bảng~\ref{tab:protocol}.
  \item Không xoá hay sửa tay file \code{outputs/preds/*.jsonl}. Script tự bỏ qua ảnh đã chạy, nên khi Colab ngắt
        chỉ cần chạy lại đúng lệnh cũ.
  \item Ghi mọi quyết định ngân sách (cắt ảnh âm, giảm epoch, mặt nạ phổi heuristic, \dots) vào
        \code{outputs/budget.md}, và nêu lại trong báo cáo.
\end{enumerate}

\paragraph{Kiểm tra bắt buộc trên Colab trước khi chạy đủ.} Làm lần lượt; mục nào không đạt thì dừng và sửa
(chi tiết trong mục ``Kiểm tra bắt buộc'' của \code{RUNBOOK.md}).
\begin{itemize}[label=$\square$]
  \item \textbf{1. Tokenizer cho điểm yes/no} (với \code{medgemma} và \code{qwen}): sau điểm phân kỳ chỉ còn chính
        từ \code{yes}/\code{no}, không có dấu nháy bị gộp; target huấn luyện sau \code{ANSWER\_PREFIX} bắt đầu bằng
        đúng token đó.
  \item \textbf{2. Một lần \code{build\_example} cho mỗi model trên ảnh thật:} không ném \code{ValueError} ở bước
        kiểm tra tiền tố prompt của Gemma.
  \item \textbf{3. Loss đầu tiên của pilot 160 mẫu} (\code{--max-train 160 --epochs 1}): \code{loss} ở bước 10
        khoảng 0,5--3. Nếu khoảng 10--25 thì gradient accumulation bị chia hai lần: dừng và sửa.
  \item \textbf{4. Loại GPU:} L4 hoặc A100, số đầu của \code{torch.cuda.get\_device\_capability()} $\ge 8$. Không
        dựa vào \code{is\_bf16\_supported()} (T4 cũng trả \code{True}).
  \item \textbf{5. Smoke YOLO 1 epoch:} file \code{best.pt} được tạo ra.
  \item \textbf{6. Thời gian mặt nạ phổi} của \code{run\_audit.py} trên 10 ảnh; nếu quá chậm dùng
        \code{--lung-mask heuristic} và ghi vào \code{budget.md}.
  \item \textbf{7. Thời gian suy luận VLM} trên 50 ảnh val (cấu hình \code{pilot}, ghi ra ngoài \code{outputs/preds/});
        ghi giây/ảnh vào \code{budget.md} và ngoại suy chi phí C0.
\end{itemize}

\begin{table}[H]
\centering
\caption{Đầu vào, lệnh, đầu ra của từng bước và chỗ trong báo cáo mà bước đó điền. Đường dẫn tương đối với thư
mục gốc dự án; \code{M = data/processed}, \code{O = outputs}.}
\label{tab:protocol}
\scriptsize
\renewcommand{\arraystretch}{1.15}
\begin{tabularx}{\textwidth}{@{}>{\raggedright\arraybackslash}p{1.05cm} >{\raggedright\arraybackslash}p{2.75cm}
  >{\raggedright\arraybackslash}p{3.15cm} >{\raggedright\arraybackslash}p{3.6cm} L@{}}
\toprule
\textbf{Bước} & \textbf{Đầu vào} & \textbf{Lệnh (RUNBOOK)} & \textbf{Đầu ra} & \textbf{Điền vào} \\
\midrule
Kiểm tra & Colab L4/A100, \code{HF\_TOKEN}, \code{M/} & 7 mục ở trên & \code{O/budget.md} & Ngân sách (Mục~\ref{sec:method}) \\
\addlinespace
E0 & Dữ liệu gốc VinBigData, NIH, JSRT & \code{prepare\_vinbig.py}, \code{prepare\_nih.py}, \code{prepare\_jsrt.py}, \code{draw\_boxes.py}, \code{data\_stats.py}
   & \code{M/manifest\_*.csv}, PNG 896, \code{O/qa/}, \code{O/e0\_data.md}, \code{O/e0\_data\_edges.json}
   & Bảng~\ref{tab:data-stats}; số ca ở Hình~\ref{fig:split} \\
\addlinespace
E1 & 3 manifest, PNG & \code{train\_densenet.py}, \code{train\_yolo.py}, seed 0--2
   & \code{O/preds/B1.jsonl}, \code{B2.jsonl}, \code{O/yolo/seed*/}\allowbreak\code{weights/best.pt}
   & Dòng B1, B2 của Bảng~\ref{tab:main}; B2 ở Bảng~\ref{tab:loc} \\
\addlinespace
E2 & 3 manifest & \code{predict\_vlm.py --config C0-medgemma|C0-qwen}, run 0 ($T{=}0$), run 1--2 ($T{=}0{,}7$), 256 token
   & \code{O/preds/C0-medgemma.jsonl}, \code{C0-qwen.jsonl}
   & Dòng C0 của Bảng~\ref{tab:main}, \ref{tab:loc}; \code{run\_agreement.csv} \\
\addlinespace
E3 & \code{manifest\_vinbig.csv} & \code{train\_lora.py --mode cls|loc} (C3: \code{--balance --pos-weight 2.0}), seed 0--2;
   \code{predict\_vlm.py --adapter}
   & \code{O/lora/C\{1,2,3\}-s*/final}, \code{timing.json}, \code{O/preds/C1|C2|C3.jsonl}, \code{<cfg>-ep<k>}
   & Dòng C1--C3 của Bảng~\ref{tab:main}, \ref{tab:loc}; mục chế độ báo bình thường \\
\addlinespace
Báo cáo lần 1 & \code{O/preds/*.jsonl}, \code{e0\_data\_edges.json} & \code{make\_report.py}
   & \code{O/report/}\allowbreak\code{val\_selection.csv}, \code{thresholds.json}, \code{summary.csv/.md}, \code{summary\_per\_seed.csv},
     \code{strata.csv}, \code{comparisons.csv}, \code{run\_agreement.csv}
   & Chọn cấu hình LoRA tốt nhất (trung bình \code{val\_sens90}); ghi vào \code{budget.md} \\
\addlinespace
E4 & Cấu hình tốt nhất từ \code{val\_selection.csv} & \code{train\_lora.py --model qwen} cùng cờ; \code{predict\_vlm.py --model qwen --config C4}
   & \code{O/lora/C4-s*}, \code{O/preds/C4.jsonl}
   & Dòng C4 của Bảng~\ref{tab:main}, \ref{tab:loc} \\
\addlinespace
E5 & $t_{90}$ từ \code{thresholds.json} (khoá \code{cfg|0|0}), adapter seed 0, \code{best.pt} seed 0
   & \code{run\_audit.py} trên \code{test} và \code{test\_jsrt}, cho C0-medgemma, LoRA tốt nhất, B2
   & \code{O/audit/}\allowbreak\code{<cfg>\_<split>.json} (\code{delta}, \code{n\_t2\_ok}, \code{flip\_t1\_all}, \code{yes\_rate\_t3}, \code{by\_size}, \code{by\_subtlety})
   & Bảng~\ref{tab:audit-res}; biểu đồ theo tầng \\
\addlinespace
Báo cáo cuối & Toàn bộ \code{O/preds/} & \code{make\_report.py} (chạy lại)
   & \code{O/report/summary.md} và các file trên
   & Mọi ô \textcolor{badred}{[--]}; Tóm tắt; Kết luận (Q1, Q2) \\
\bottomrule
\end{tabularx}
\end{table}

\paragraph{Lịch dự kiến.} Hình~\ref{fig:gantt} là lịch đề xuất cho phần việc còn lại. Ngày tháng chỉ là đề xuất,
nhóm điều chỉnh theo thời gian đo được ở kiểm tra mục 3 và 7.

\begin{figure}[H]
\centering
\begin{tikzpicture}[x=0.32cm,y=1cm]
  % trục: ngày 0 = 14/11/2026, ngày 34 = 18/12/2026
  \foreach \d/\l in {0/14/11,7/21/11,14/28/11,21/05/12,28/12/12,34/18/12}{
    \draw[black!15] (\d,0.35) -- (\d,-5.25);
    \node[fnote, anchor=south] at (\d,0.35) {\l};}
  \newcommand{\gtask}[5]{% hàng, bắt đầu, kết thúc, màu, nhãn
    \node[font=\scriptsize, anchor=east] at (-0.4,-#1*0.55) {#5};
    \fill[#4!35, draw=#4] (#2,-#1*0.55-0.19) rectangle (#3,-#1*0.55+0.19);}
  \gtask{0}{0}{4}{black}{Kiểm tra trước + E0}
  \gtask{1}{4}{10}{okgreen}{E1 baseline B1, B2}
  \gtask{2}{5}{11}{black}{E2 zero-shot C0}
  \gtask{3}{7}{19}{uitblue}{E3 LoRA C1--C3}
  \gtask{4}{19}{20.5}{warnorange}{\code{make\_report} + chọn cấu hình}
  \gtask{5}{20.5}{25}{uitblue}{E4 C4 (Qwen)}
  \gtask{6}{20.5}{26}{badred}{E5 kiểm toán}
  \gtask{7}{26}{27.5}{warnorange}{\code{make\_report} lần cuối}
  \gtask{8}{14}{31}{black}{Viết báo cáo}
  \gtask{9}{28}{33}{black}{Slide, rà soát}
  \fill[warnorange] (20.5,-4*0.55) circle (3pt);
  \node[fnote, anchor=west, text=warnorange] at (21.2,-4*0.55) {cổng chọn (chỉ validation)};
\end{tikzpicture}
\caption{Lịch \textbf{đề xuất} cho phần việc còn lại, từ 14/11 đến 18/12/2026; chưa phải lịch đã chốt. Màu theo
loại việc như các hình trước: xanh lá baseline, xám zero-shot, xanh dương LoRA, cam bước chọn và tổng hợp, đỏ kiểm toán.}
\label{fig:gantt}
\end{figure}

\subsection{Prompt và định dạng đầu ra}

\begin{lstlisting}[escapeinside={(*}{*)}]
(*\lstcmt{\# Câu hỏi định vị (C0, C2--C4)}*)
Is there a lung nodule or mass on this chest X-ray? If yes, give the bounding
box(es) in pixel coordinates. Answer in JSON.

(*\lstcmt{\# Câu hỏi chỉ phân loại (C1)}*)
Is there a lung nodule or mass on this chest X-ray? Answer in JSON.

(*\lstcmt{\# Đáp án đích (C0, C2--C4); C1 chỉ có khoá nodule}*)
{"nodule": "yes", "boxes": [[412, 230, 468, 281]]}   (*\lstcmt{\# ca dương, toạ độ trên ảnh 896$\times$896}*)
{"nodule": "no", "boxes": []}                        (*\lstcmt{\# ca âm}*)
\end{lstlisting}

Parser chấp nhận JSON hợp lệ có khoá \code{nodule} nhận \code{yes}/\code{no}. Box không hợp lệ (toạ độ ngoài
ảnh, $x_2 \le x_1$) bị loại. Câu trả lời không parse được tính là ``no'' khi tính sensitivity/specificity thô,
bị loại khỏi các chỉ số chế độ báo bình thường, và tỷ lệ của chúng được báo riêng.

\subsection{Siêu tham số}

\begin{table}[H]
\centering
\caption{Siêu tham số huấn luyện.}
\label{tab:hparams}
\small
\begin{tabularx}{\textwidth}{@{}p{4.2cm} L@{}}
\toprule
\textbf{Thành phần} & \textbf{Giá trị} \\
\midrule
LoRA & $r = 16$, $\alpha = 32$, dropout 0,05; attention + MLP của LLM \\
Bộ mã hoá ảnh & đông cứng \\
Tối ưu & AdamW, lr $10^{-4}$, cosine, warmup 3\% \\
Epoch, batch & 3 epoch; batch 1 $\times$ tích luỹ gradient 8 \\
Độ chính xác & bf16, gradient checkpointing \\
Ảnh âm/ảnh dương khi huấn luyện & tối đa 4:1 \\
C3 & oversampling 1:1 (\code{--balance}); $w_{+} = 2$ (\code{--pos-weight 2.0}), áp lên mọi token đáp án ca dương \\
Phần cứng & GPU compute capability $\ge 8$ (L4/A100) cho MedGemma \\
Sinh văn bản & tối đa 256 token ở C0; 96 token (mặc định) ở C2--C4, 16 token ở C1; nhiệt độ 0 (thêm 2 lần 0,7 ở E2) \\
B1 DenseNet-121 & Adam, lr $10^{-4}$, 10 epoch, batch 32, BCE có \code{pos\_weight} \\
B2 YOLOv8 & ảnh 896, 50 epoch, batch 16 \\
\bottomrule
\end{tabularx}
\end{table}

\subsection{Mã nguồn và tái lập}

Mã nguồn nằm trong gói Python \code{nodulecxr} (thư mục \code{code/}), viết theo hướng kiểm thử trước, gồm các
module: chuẩn bị dữ liệu và hợp nhất box, prompt và parser, huấn luyện LoRA, baseline CNN, chỉ số phát hiện và
định vị, thống kê, kiểm toán che vùng, và tổng hợp bảng. Các bước cần GPU chạy trên Colab theo \code{RUNBOOK.md}.
Bảng kết quả trong báo cáo được sinh từ file dự đoán JSONL, không gõ tay.

\subsection{Mức kiểm chứng tài liệu}

Tổng quan dựa trên 90 bài (\code{docs/lit-review/lung-cancer-cxr/corpus.md}). Hai bài cốt
lõi~\cite{khovanova2026comparative,nishio2026zeroshot} đã qua bình duyệt và được đọc toàn văn
(\cite{nishio2026zeroshot} qua bản preprint medRxiv). Các bài preprint
(\cite{muthyala2026frozen,choi2026chexpercept,fisher2025vlfm,mustafic2026noduloc}) chưa qua bình duyệt và được
trích dẫn thận trọng.
